\documentclass[b5paper]{report}
\usepackage[lmargin=25mm,rmargin=25mm,tmargin=27mm,bmargin=30mm]{geometry}
\usepackage{color}
\usepackage{xcolor}
\usepackage{tikz}
\usepackage{pgfplots}
\usepackage{graphicx}
\usepackage{wrapfig}
\usepackage{svg}
\usepackage{float}
\usepackage{url}
\usepackage{parskip}
\usepackage{hyperref}
\usepackage[utf8]{inputenc}
\usepackage[chapter]{minted}

\graphicspath{{images/}}
\svgpath{{images/}}
\usetikzlibrary{positioning}
\usetikzlibrary{shapes,decorations,shadows}

\begin{document}

\begin{titlepage}
  \centering
	\includegraphics[width=0.55\textwidth]{ntnu}\par\vspace{1cm}
	{\scshape\Large TDT4900 Master's Thesis \par}
	\vspace{1.5cm}
	{\huge\bfseries Durability in a data-flow based storage system\par}
	\vspace{2cm}
	{\Large\itshape Lars Martin Bævre Ek \\ \texttt{larek@stud.ntnu.no} \par}
	\vfill
	supervised by\par
  Professor Kjetil \textsc{Nørvåg}

	{\large \today\par}
\end{titlepage}

\newcommand{\note}[1]{{\color{olive}\sffamily #1}}
\newcommand{\todo}[1]{{\color{purple}\sffamily[TODO: #1]}}
\newcommand{\eg}{\textit{e.g.},\xspace}
\newcommand{\ie}{\textit{i.e.},\xspace}
\newcommand{\vs}{\textit{vs.}\xspace}
\newcommand{\viz}{\textit{viz.},\xspace}

% Handy superscript macros
\newcommand{\stss}{\textsuperscript{st}}
\newcommand{\ndss}{\textsuperscript{nd}}
\newcommand{\rdss}{\textsuperscript{rd}}
\newcommand{\thss}{\textsuperscript{th}}

% Enumeration commands
\newcommand{\one}{({\em i}\/)}
\newcommand{\two}{({\em ii}\/)}
\newcommand{\three}{({\em iii}\/)}
\newcommand{\four}{({\em iv}\/)}
\newcommand{\five}{({\em v}\/)}

\newcommand\code[1]{\texttt{#1}}


\begin{abstract}
\end{abstract}

\tableofcontents
\pagebreak

\chapter{Introduction}

\chapter{Background} \label{chap:background}

\section{RocksDB}
RocksDB is an embedded key-value store optimized for modern flash storage.
RocksDB started out at Facebook, with the goal of making a version of Google's
LevelDB that performed well on modern hardware. Today, RocksDB is used at the
heart of a wide variety of databases, such as CockroachDB, MyRocks (Facebook's
fork of MySQL) and TiDB.

Traditional B-tree based database systems are often faced with poor write
performance as a result of random writes, which perform worse than sequential
writes on both magnetic and flash based storage mediums. RocksDB, on the other,
hand achieves impressive write performance through the use of immutable
log-structured merge trees (LSM-trees), avoiding the need for random writes to
persistent storage altogether.

Writes are initially only written to a persistent write-ahead log (WAL) and
in-memory data structures referred to as memtables. Later on these memtables are
flushed to RocksDB's LSM-tree data structure, Static Sorted Tables (SST). The
latter is done by background threads, allowing regular processing to continue
without getting blocked by slow writes to persistent storage.

% TODO:
% * Calling into RocksDB (C-API, rust-rocksdb)
% * Iterators
% * Prefix Iteration/Slice Transform

\subsection{MemTables}
All writes are initially synchronously written to an in-memory data structure--a
memtable--which is later flushed to disk at the point of filling up. Both the
size and the number of memtables can be configured at runtime.

RocksDB's default memtable implementation is a skiplist, with an $ O(log n) $
bound on inserts, searches and deletes. This can be changed to a series of hash
based implementations, which offer better performance if all operations are done
within a pre-specified key prefix.

\subsection{Static Sorted Tables}
After a memtable reaches a certain size, RocksDB's background threads will be
responsible for flushing it to persistent storage. This will generate one or
more SS-tables on disk, where each file is sorted. SS-tables are immutable: a
new SS-table is always created, and existing ones are never updated. This
ensures that writes remain sequential.

\subsection{Basic Operations}
Akin to other key-value databases, RocksDB offers a familiar API of
\code{Put(key, value)}, \code{Get(key)} and \code{Delete(key)}, operating
directly on bytestream values. Both insertions and deletions are purely
sequential: subsequent \code{Put} operations of the same key never backtrack and
overwrite existing keys, and deletions insert tombstone markers to avoid having
to randomly read and mutate previously written values.

This means that a read operation has to search through all the potential row
destinations to find its target value.
target value.


This is part of RocksDB's
promise of write performance, and naturally leads to slightly worse read
performance compared to something like a B-tree.



\subsection{Compactions}
To ensure purely sequential writes, keys are never overwritten directly by
RocksDB.

SS-tables are organized in levels, where a new level is created when the
previous one reaches a certain size.

This is done in the background by RocksDB's
worker threads




RocksDB maintains writes in an in-memory buffer data structure

, built
by Facebook. The project started out as a fork of of Google's LevelDB--another
key-value store--with the goal of building a database targeting modern hardware.

RocksDB achieves impressive write performance through the use of log-structured
merge trees.


\bibliographystyle{abbrv}
\bibliography{sources}

\end{document}
